\section{Finite onset conditions and order-dependent constants}
\label{app:onsets}
The following sufficient conditions make the fixed-order theorem
fully parameterized. In addition to the classical real
Riemann--Siegel and sectorial gamma onsets, impose
\begin{subequations}\label{eq:structural-gates}
\begin{align}
 &L\ge16,\quad X\ge\max(1,4(R+Y+1)),\quad V_B\le X/4,\quad
                         X-V_B\ge X/2,\\
 &a(X+V_B)-a(X-V_B)<1,\qquad
                       2p\le (X-V_B)/(4\pi),\\
 &DV_B/X\le1/2,\qquad 2AV_B^2/X\le1,\\
 &D_0/X\le1-\pi/8,\qquad
 B\{(D_0/X)(\pi/4+L/2)+D_0^2/X^2\}\le1,\\
 &V_B\ge\max(2R,W+1,32B).
\end{align}
\end{subequations}
Here \(D_0\) is a fixed constant in the sectorial slope formula
\[
 \lambda_{\epsilon n}=-\pi/8+\epsilon\ii v_n/2+\delta_\epsilon,
 \qquad
 v_n=\log(a_X/n),\qquad \abs{\delta_\epsilon}\le D_0/X.
\]
The optional inherited real-chart restriction
\(V_B\le\log^2X\) is also eventual.
Writing \(\log^+a=\max(0,\log a)\), the fixed-order tail
conditions are
\begin{subequations}\label{eq:fixed-tail-gates}
\begin{align}
 L^2&\ge(P_0+p-1)\log X+\log^+ A_{\mathrm{mod}},\\
 L^2&\ge(p+1)\log X+\log^+ A_{\mathrm{art},p},\\
 X^2/(128B)&\ge\pi X/8+(p-1)\log X+\log^+ A_{\mathrm{far}}.
\end{align}
\end{subequations}
They put each tail in \cref{sec:comparison} below
\(X^{-(p+3/4)}\). Every condition is eventual for fixed
\(p,\tau,B,r,b\). The factors
\(\e^{Y^2/(4\tau)}\), \(\tau^{-1}\), \(\tau^{-2}\),
and the full moments in \cref{eq:gaussian-budgets,eq:tail-moments}
remain in these constants.

\subsection{Common gates at growing order}
\label{app:growing-gates}
For \(1\le p\le cL\), \(B=2,\tau=1\), replace the retained-order
condition in \eqref{eq:structural-gates} by
\(2cL\le(X-V)/(4\pi)\).
Include the fixed outer-calibration gate
\(C_{\mathrm{cal}}(1+L)/X\le1\).
Do not impose \eqref{eq:fixed-tail-gates} at growing order.
Instead, with \(C_{\mathrm o}=C_\chi/c_{\mathrm o}\), the following
four inequalities suffice for \cref{thm:growing}:
\begin{subequations}\label{eq:growing-gates}
\begin{align}
 (p+3/4)\log X
 &\ge(\eta+1/8)L^2+\log^+(4C_{\mathrm o}Q_p),\\
 (16\theta-\eta-1/8)L^2
 &\ge(P_0-7/4)\log X+
                   \log^+(196C_{\mathrm o}A_{\mathrm{mod}}^\theta),\\
 (16\theta-\eta-1/8)L^2
 &\ge\tfrac14\log X+
                   \log^+(196C_{\mathrm o}A_{\mathrm{art},p}^\theta),\\
 X^2/256
 &\ge\pi X/8+(\eta+1/8)L^2+
                   \log^+(196C_{\mathrm o}A_{\mathrm{far}}).
\end{align}
\end{subequations}
Each pays at most one quarter of the target.
For the half-target truncation required in
\cref{sec:evaluation}, replace \(4,196\) by \(8,392\).
These are the operative finite inequalities.
\Cref{thm:explicit-onset} now supplies an explicit sufficient all-height
threshold; a smaller height can be accepted by its finite verifier.
For the moving domain use the distinct constants
\eqref{eq:moving-truncation-constants}, unchanged half-target factors
\(8,392\), and the continuous slack in
\cref{subsec:moving-onset-proof}. Do not substitute the moving endpoint
into the fixed-domain constants above.

\section{Evolved population bounds and finite level size}
\label{app:population}
The error proof used the real initial population. One can also
bound the entire evolved population without discarding the
near-cutoff modes. Put \(a_0=\pi/8\), \(c_Y=(Y-1)/2\),
\(c_Y^+=\max(c_Y,0)\), and for an integer \(d_0\ge0\) define
\begin{subequations}\label{eq:Pi}
\begin{align}
 W_{d_0}(\tau)
 &=2^{d_0}+\frac{3^{d_0}}2\sum_{k=0}^{d_0}\binom{d_0}{k}
              (8/\tau)^{(k+1)/2}\Gamma((k+1)/2),\\
 \Pi_{d_0}(\tau,B)
 &=4(4\pi)^{-1/4}4^{d_0}
   \exp\{R+Y+Ba_0^2+1+(Y+1)/2+2(c_Y^+)^2/\tau\}
                   W_{d_0}(\tau).
\end{align}
\end{subequations}
The slope gates in \eqref{eq:structural-gates} imply
\[
 \abs{\e^{T\lambda_{\epsilon n}^2}}
       \le\e^{Ba_0^2+1}\e^{-\tau v_n^2/4},
 \qquad
 \abs{\e^{\lambda_{\epsilon n}d}}
       \le\e^{R+Y}\e^{Yv_n/2}.
\]
Also \(1+\abs\lambda\le4(1+v_n)\).
The shell \(j\le v_n<j+1\) contains at most \(2a_X\e^{-j}\)
integers. This remains true in the last nonempty shell,
where \(a_X\e^{-j}\ge1\).
In that shell \(n^{-1/2}\le a_X^{-1/2}\e^{(j+1)/2}\).
Use
\[
 c_Yj-\tau j^2/4
       \le2(c_Y^+)^2/\tau-\tau j^2/8.
\]
For \(j\ge1\), compare with the integral on \([j-1,j]\);
\(j+2\le3(v+1)\) there and
\(\e^{-\tau j^2/8}\le\e^{-\tau v^2/8}\).
The half-Gaussian moments give exactly the bound \(W_{d_0}\).
Summing both arms proves
\begin{equation}
 \sum_{\epsilon,n\le N}
 \abs{b_{\epsilon n}\e^{\lambda_{\epsilon n}d+T\lambda_{\epsilon n}^2}}
       (1+\abs{\lambda_{\epsilon n}})^{d_0}
       \le \Pi_{d_0}(\tau,B)gX^{1/4}.
 \label{eq:evolved-population}
\end{equation}

For the moment polynomial define
\[
 Z_m(B)=m!\sum_{\nu=0}^{\lfloor m/2\rfloor}
       \frac{B^\nu(W+2B)^{m-2\nu}}{\nu!(m-2\nu)!}.
\]
Since \(\abs{d+2T\lambda}\le(W+2B)(1+\abs\lambda)\),
the whole \(m\)-th moment sum is at most
\(Z_m\Pi_m gX^{1/4}\).
For \(j\ge1\) this gives
\begin{equation}
 \sum_{\epsilon=\pm1}\abs{U_{\epsilon,j}}
 \le gX^{1/4-j}
       \sum_{k=1}^j\frac{A^kD^{j-k}}{k!}\binom{j-1}{k-1}
                                   Z_{j+k}\Pi_{j+k}.
 \label{eq:level-size}
\end{equation}
For \(\tau=1/L\), the degree-dependent comparison
\eqref{eq:moving-evolved-ratio} is the separate evolved-population
ledger; the order-independent truncation multiplier
\eqref{eq:moving-truncation-ratio} does not bound it.
On the slab \([1,2]\), the explicit factorial sums give fixed
\(K_0,K_1\) such that \(Z_m\Pi_m\le K_0K_1^m(m!)^2\).
For example, \(Z_m/m!\) is a coefficient of
\(\exp((W+4)s+2s^2)\), so \(Z_m\le m!\e^{W+6}\);
the half-gamma bound similarly controls \(\Pi_m\).
Consequently the \(j\)-th level has the absolute envelope
\[
 gX^{1/4-j}\exp(O((j+1)\log(j+2))).
\]
This is a bound for the entire population, not a signed leading
asymptotic, a monotonicity statement, or a convergence theorem
for the sum over all levels.

\section{Sources, scope, and manuscript correspondence}
\label{app:provenance}
Internal note paths in the bibliography are relative to
\nolinkurl{papers/investigations-5/01_heat_density_asymptotics/}.
The analytical inputs and their derived uses are deliberately
separated:
\begin{center}
\small
\begin{tabular}{@{}>{\raggedright\arraybackslash}p{.26\textwidth}>{\raggedright\arraybackslash}p{.34\textwidth}>{\raggedright\arraybackslash}p{.29\textwidth}@{}}
\toprule
Input or research note & Use here & Scope retained\\
\midrule
Polymath Proposition 6.2 &
Real canonical remainder for each positive \(K\) &
No evolved negative-time theorem imported\\[3pt]
DLMF gamma asymptotics &
Sectorial derivative and calibration bounds &
Fixed sectors, explicit eventual gates\\[3pt]
Arias exact identities &
Weighted coefficients and \(F_0\) derivatives &
Preliminary precision conclusions excluded\\[3pt]
DLMF Gauss properties &
Positivity and degree \(2n-1\) exactness &
Complex error bound derived here\\[3pt]
Recurrence Notes 02--04 &
Finite hierarchy and compact comparison &
Fixed \(p,\tau,B\) first\\[3pt]
Recurrence Note 05 &
Simultaneous order costs and outer accuracy &
\(p=O(L)\) on \([1,2]\)\\[3pt]
Recurrence Note 06 &
Regularity, conditioning, and evaluation budget &
Modular provider and arithmetic contracts\\[3pt]
Recurrence Note 09 &
Finite elementary, scale, coefficient and guard constructions &
Strict descriptors and exact queries\\[3pt]
Recurrence Note 10 &
Shifted gamma, Gauss root counts, total bit work and storage &
All modes and actual separation retained\\[3pt]
Recurrence Note 11 &
Uniform constants and explicit all-height onset &
Fixed radii, \(1\le T\le2\)\\[3pt]
Recurrence Note 12 &
Linked finite component replay &
Supporting evidence only\\[3pt]
Recurrence Note 16 &
Additive-cost linear sufficient log-height &
All original gates and half-target budget\\[3pt]
Recurrence Note 18 &
Sharper uniform RS bounds and contour count &
Same domains and delivered accuracy\\[3pt]
Recurrence Note 20 &
Integrated moving-endpoint proof in \cref{sec:moving-endpoint} &
\(L^{-1}\le T\le2\), strict queries and paid tag\\[3pt]
Book manuscript &
Common notation and cross-region synthesis &
Proof ownership remains here\\[3pt]
Ansatz Notes 17--20 and companion &
Positive finite hierarchy and compact bridge &
Fixed positive order and its own paid domains\\
\bottomrule
\end{tabular}
\end{center}
The preceding ansatz results in \cref{subsec:retention} are
retained as separately sourced theorems with their original
domains. Their first-correction formula is reconciled with
the present one in \cref{subsec:low-orders}; no differentiation
rule is transferred from one construction to the other.
For the current positive comparison use
\cref{subsec:positive-comparison} and the proof labels in
\cite{ansatzcomp}; the Book's negative-time contracts point to
\cref{thm:fixed,thm:growing,thm:linear-onset,thm:effective,thm:bit-cost},
\cref{thm:moving-endpoint-delivery} and
\cref{prop:sharper-rs,prop:smaller-contour} here.
Note 19 retains AN08's original request. The later Note 20/Review 14
return is Complete/CLOSED in the live register; its full proof is now
integrated in \cref{sec:moving-endpoint}
\cite{rec19,rec20,rec20review,recstatus}. The input interface and
provider cost changes are given in \cref{subsec:moving-input,subsec:moving-cost}.
The repository result index and control screen
\cite{results,controls}, together with its scoped-error cautions
\cite{avoid}, govern the interpretation.

The source notes were developed and internally assessed with
gpt-6.1-sol, recorded effort ultra. This manuscript was drafted
for Edward Baker with gpt-6-astra, recorded effort max;
the linear-onset and uniform RS-bound strengthenings use gpt-6.1-sol,
actual effort ultra. The Book-companion reconciliation and subsequent moving-endpoint proof
integration were prepared with gpt-6.1-sol, actual effort ultra.
The mathematical status remains that of the stated proofs
from named inputs. The modular contracts are supplied by finite
constructions on the declared input class; complete software
certification and independent specialist review remain open.
Compilation, reference checks, and visual inspection concern
document production. They are not numerical evidence for
the analytical estimates.
